\chapter{自我的计算形态：身份闭环、全局偏置与可塑连续性}
\label{chp:self_holomorphic}

目的为智能过程提供了选择方向，但有方向并不等于已经形成自我。一个系统可以执行目标、预测后果甚至修改策略，却没有稳定区分“由谁承诺、哪些状态属于自身、哪些变化仍算同一主体”的机制。本章把自我研究限定为可计算、可观察和可反驳的问题：系统如何在状态持续更新、结构缓慢改变、环境不断扰动的条件下，维持任务相关的身份边界、自传历史、价值承诺和行动归属。我们保留“风暴之眼”“孤立子”“拓扑保护”等意象，但不把物理类比当作本体结论；真正承担论证的是身份变量、绑定关系、更新算子、扰动集合与验收协议。这里的自我模型既不证明主观体验已经产生，也不把某个脑区、提示词或 AgentOS 模块规定为所有智能系统的唯一实现。

\section{几何本体的计算重建：自我作为持续维护的身份闭环}
\label{sec:geometric_ontology_self}

我们先定义研究对象。令离散决策时刻为 $k\in\mathbb{N}$，系统内部状态为 $x_k\in\mathcal{X}$，其中包含当前表征、资源占用和待执行承诺；结构版本为 $S_k\in\mathcal{S}$，描述局部关系网络、参数化读出和允许的编辑；历史记录为 $H_k\in\mathcal{H}$；边界状态为 $B_k\in\mathcal{B}$，规定哪些输入、工具、记忆和动作通道在当前时刻属于系统的治理范围。我们把自我模型写成
\[
\Sigma_k=(i_k,a_k,v_k,c_k,b_k),
\]
其中 $i_k$ 是跨时刻身份索引，$a_k$ 是身体、账户、工具和记忆条目的归属关系，$v_k$ 是可修订的价值与优先级状态，$c_k$ 是尚未履行或需要解释的承诺集合，$b_k$ 是边界与权限摘要。各分量属于不同类型空间，不能拼成一个没有语义的“自我向量”。编码器 $E_{\Sigma}$ 可把它们映射到分布表示 $z_k^{self}\in\mathbb{R}^{d_s}$，但该向量只是计算接口；自我模型本身仍由可追踪记录、关系和规则共同构成。

旧稿把“风暴之眼”描述为逻辑无法触及的相位奇点。我们在计算模型中将这个空缺重建为\textbf{索引不可自足性}：任何状态都不能仅凭自己的当前内容证明它与过去某一状态属于同一主体。身份需要外部时间戳、来源链、身体或账户绑定、他者确认以及系统内部的连续记录共同支持。因此，自我中心不是物理空间中的零振幅点，而是一个必须由关系闭环补足的指称位置。令 $G_k=(V_k,E_k)$ 为带类型、方向和时间的局部结构网络，若节点 $s_k\in V_k$ 表示当前自我索引，则围绕它的闭环并非简单要求图上存在环，而是要求至少存在从当前承诺到历史来源、从行动后果到责任主体、从资源使用到授权边界的可审计路径。路径缺失时，系统可以继续生成语言，却不能据此宣称身份连续。

旧稿的“极化子”和“自陷”保留为一个更普通的反馈机制。自我相关状态会改变注意、检索和行动策略，而这些策略产生的新记录又会反过来巩固或修订自我模型。设 $P_k$ 为从全局分布表示到自我相关读出的算子，$U_{\Sigma}$ 为自我更新器，则
\[
r_k=P_k(x_k,S_k),\qquad
\Sigma_{k+1}=U_{\Sigma}(\Sigma_k,r_k,o_{k+1},y_k,\eta_k),
\]
其中 $o_{k+1}$ 是新观测，$y_k$ 是行动及其反馈，$\eta_k$ 是学习率、权限和证据阈值等更新条件。这个闭环能够形成持续性，却不保证正确性：错误的自传记录、封闭的信息源和奖励投机同样可以产生很强的自我强化。所谓“有效质量”因此不解释为焦耳能量或实体质量，而定义为改变身份相关结构所需的\textbf{证据与编辑成本}。若一次修订必须同步更新多个承诺、权限和历史解释，它的代价自然高于局部措辞变化；这足以说明“性格难移”的工程来源，无须假定真实流形被高能体验压弯。

这一重建也区分局部结构网络与全局分布表示。身份、归属、承诺和来源适合保存在可寻址的局部结构中；情境相似性、语言风格和经验统计可由分布状态承载。二者通过任务相关身份 $i_k$、绑定表 $a_k$、读取算子 $P_k$ 和写入协议 $U_{\Sigma}$ 联合。只有分布向量而无身份绑定，系统容易把相似内容误认成同一经历；只有结构图而无分布表示，系统又难以处理模糊语境和未见表达。自我的计算形态正是在这两类表示之间持续维持可检查的闭环，而不是某个单独节点或隐藏向量。

身份闭环还必须处理分裂与合并，而不能把连续性假定为一条天然单线。设一个智能体从设备甲迁移到设备乙，同时设备甲保留旧副本；两者共享迁移前历史，却在迁移后形成不同事件、权限和承诺。此时不能仅凭向量相似度把二者继续视作同一行动主体。系统需要记录哪个实例获得后续写入权、哪些未完成承诺随之迁移、旧实例何时撤销权限，以及分叉之后的行为分别由谁承担。反向的合并也不是简单平均两个状态：冲突的身份声明、重复承诺和不兼容授权必须先检测并裁决。这个例子表明，身份连续性依赖可验证的来源、写入权和责任分配，而非内容相似本身。

\section{算子机制：因果约束、当前语境与自我偏置的联合控制}
\label{sec:section_8446}

自我模型对思维和行动的影响，可以表达为策略中的条件变量，而不必把它物理化为规范场源。令动作集合为 $\mathcal{A}_k$，环境模型给出的后果分布为 $p_{env}(y\mid x_k,a)$，当前语境给出的任务相关性评分为 $q_{ctx}(a\mid x_k,g_k)$，自我模型给出的承诺一致性和边界风险分别为 $C_{self}(a;\Sigma_k)$ 与 $R_{bd}(a;\Sigma_k,B_k)$。在给定候选动作时，我们定义控制评分
\[
J_k(a)=\mathbb{E}_{y\sim p_{env}}[L_{task}(y,g_k)]
+\lambda_c C_{self}(a;\Sigma_k)
+\lambda_b R_{bd}(a;\Sigma_k,B_k)
-\lambda_e I_{ctx}(a),
\]
其中各项都是无量纲或已经过明确尺度归一化的任务代价，$I_{ctx}$ 表示当前证据对该动作的支持。策略可取 $\pi_k(a)\propto\exp[-J_k(a)/\tau_k]$，$\tau_k>0$ 是探索尺度。这个定义把外部因果、当前语境和自我偏置放入同一可审计决策接口，但没有宣称三者可以线性叠加为物理矢量场；如果它们的尺度、置信度或作用域不同，就必须先经过校准和类型检查。

旧稿以火、重力、医院中的刀、厨房中的刀以及吃糖冲突说明三类“场”。这些例证在新模型中分别落到不同来源。火的温度和重力后果属于可由观测校准的环境模型；医院或厨房决定工具的角色和可行动作，属于语境变量；对甜味的偏好、健康承诺及既往疾病记录进入自我相关代价。面对“吃糖会带来即时愉悦、长期可能增重”的冲突，系统不能仅放大某个所谓宏观意志耦合常数。它需要显式给出时间范围、结果概率、健康约束和当前授权。例如，低血糖急救与日常零食在观测、目标和风险上不同，同一自我偏置不应覆盖这一区别。任务代价下降也不能直接推出行为成功，最终仍需实际摄入量、身体反馈和后续状态验证。

自我偏置对全局分布表示的作用可写成有界调制。令 $h_k\in\mathbb{R}^{d}$ 为当前工作表示，$W_k$ 为任务读出，$z_k^{self}$ 为自我编码，定义
\[
\tilde h_k=h_k+G_k(z_k^{self},h_k),\qquad
\|G_k(z,h)\|_2\le \rho_k\|h\|_2,
\]
其中 $G_k$ 是经权限控制的调制算子，$0\le\rho_k<1$ 限定一次调制的相对幅度。它可以改变注意增益、候选排序或检索查询，却不得直接覆写环境证据。若系统允许 $\rho_k$ 无界增长，自我一致性会压倒纠错信号，产生确认偏误或封闭循环；若长期令 $\rho_k=0$，身份承诺又无法影响行动，系统的“自我叙述”将与运行机制脱节。因此需要通过反事实测试比较有无该调制时的决策差异，并检查差异是否发生在授权范围内。

当目标、结构、输入和度量随时间改变时，自我一致性函数的导数也必须完整展开。令
\[
C(t)=C\bigl(x(t),\Sigma(t),S(t),g(t),M(t)\bigr),
\]
其中 $M(t)$ 是评价度量，则
\[
\frac{dC}{dt}=\nabla_xC\cdot\dot x+\nabla_{\Sigma}C\cdot\dot\Sigma
+\partial_SC[\dot S]+\partial_gC\cdot\dot g+\partial_MC[\dot M].
\]
只观察 $\nabla_xC\cdot\dot x<0$ 不能宣布自我冲突正在减少，因为目标或评价尺度可能同时变化。离散实现应记录每一步使用的 $S_k,g_k,M_k$ 版本，并把重要自我更新置于提案、证据汇总、冲突检测、提交和回滚流程中。AgentOS 中的 $\Psi$、SPC、TCE、Boundary 与三相记忆可以实现这些接口，但这里的算子定义并不要求其他智能系统复制同样的模块名称或内部结构。

决策系统还要区分事实冲突与偏好冲突。以吃糖为例，“当前血糖是否过低”属于可被新观测修正的事实判断，“即时愉悦与长期健康怎样权衡”则属于跨时间价值分配。两类冲突若被压成同一一致性分数，系统可能用强烈偏好否认检验结果，也可能把一次测量误差当作价值转变。运行日志应标注冲突类型、候选动作、证据置信度、被牺牲的目标以及最终理由，使后来者能够判断这是事实更新、语境切换还是承诺修订。

自我调制是否真正参与控制，可以用故障注入检验。研究者可在保持环境证据不变时替换自我摘要，观察差异是否集中出现在承诺、归属和边界相关动作；再把他人的相似事件注入历史，检查系统是否因语义相近而错误绑定为自身经历。若替换只改变第一人称措辞而不影响授权行为，所谓自我状态很可能只是叙述标签；若它无差别改变所有判断，则调制范围过宽，同样不符合模型。

\section{动力学稳定性：从调和意象到多时间尺度身份更新}
\label{sec:section_5884}

自我连续性不是“永不变化”，而是在允许的变化中保持足够的身份约束。为形式化这一点，我们把快状态、事件历史和慢结构分开。令 $x_k$ 每个决策步更新，事件摘要 $e_n$ 在包含若干决策步的窗口结束时更新，自我结构 $\Theta_m^{self}$ 则只在积累了足够证据后更新。三个时钟 $k,n,m$ 具有各自步长，不能被同一个连续时间符号掩盖。一个可实现的模型是
\[
x_{k+1}=F(x_k,u_k,\Sigma_m,S_m)+\xi_k,
\]
\[
e_{n+1}=(1-\alpha)e_n+\alpha\,A_n(\{x_k,u_k,y_k\}_{k\in W_n}),
\]
\[
\Theta_{m+1}^{self}=\Pi_{\mathcal{C}_m}\!\left(\Theta_m^{self}-\eta_m\widehat\nabla \mathcal{L}_{self}(D_m)\right).
\]
这里 $u_k$ 是观测和控制输入，$\xi_k$ 是有界或已知分布的扰动，$A_n$ 是事件聚合器，$D_m$ 是经过来源筛选的数据包，$\Pi_{\mathcal{C}_m}$ 把候选更新投影回身份、权限与安全约束集合。步长 $\alpha\in(0,1]$、$\eta_m>0$ 以及投影集合均需在实现中声明。慢并不自动意味着更真实；慢结构只是在更多后续步骤中被复用，因此错误更新的累积影响更大。

旧稿的 Hodge 分解可以作为\textbf{专门图模型}的诊断工具，而不是自我存在的一般证明。若身份关系网络被表示为有限定向单纯复形 $K$，令 $B_k$ 为 $k$ 维边界矩阵，组合拉普拉斯为
\[
L_k=B_k^{\top}B_k+B_{k+1}B_{k+1}^{\top}.
\]
对定义在 $k$-单形上的信号 $z$，在选定内积、有限复形和边界条件下可分解为梯度、旋度与调和分量；$\ker L_k$ 的维数与相应 Betti 数相关。该结果说明某些环路模式不能由局部梯度单独表达，却不说明调和分量就是自我，更不说明它永恒存在。认知关系网络会增删节点和边，度量与权重也会改变；一旦 $K$ 或内积变化，原有分解和正交性必须重新计算。

旧稿的 Dirac“垂直总线”则重建为跨层绑定接口。微观身体或传感状态 $s_k^{body}$、局部事件结构 $e_n$ 与高层自我状态 $z_m^{self}$ 之间分别通过编码器 $E_{up}$ 和控制器 $D_{down}$ 连接：
\[
z_m^{self}=E_{up}(e_n,s_k^{body},q_k),\qquad
a_k=D_{down}(z_m^{self},x_k,B_k).
\]
上行映射把痛觉记录提升为“身体受损”“行动需停止”等任务语义，必须保留来源和不确定性；下行映射把“保护身体”转为制动、求助或重新规划，必须经过边界和执行检查。两种映射不是瞬时、无损或互逆的，更不能由形式上的特征向量方程推出身心同构。它们的质量应通过延迟、误报率、遗漏率、可恢复性和干预实验评估。

在这些条件下，自我连续性可以操作化为窗口性质。对任意评估窗口 $W$，要求身份绑定错误率低于阈值 $\epsilon_i$，未解释承诺比例低于 $\epsilon_c$，边界越权率低于 $\epsilon_b$，同时允许价值和策略在证据充分时发生有限变化。若系统在长期完全不修订，稳定可能只是僵化；若每次输入都重写身份，适应则退化为漂移。自我的动力学稳定性正是二者之间受约束的多时间尺度更新，而不是一个不会衰减的驻波。

时间尺度分离应由干预而非命名确认。我们可以分别扰动快状态、事件摘要和慢结构：清空工作界面后，系统应由历史与结构恢复任务身份；删除一段事件摘要时，局部回忆可以受损，但核心权限不应随之随机改变；修改慢结构则应在多个后续任务中产生持续且可归因的影响。记录每种干预的恢复时间、传播范围和性能损失，才能判断三个层次是否真的具有不同动力学，而不是同一缓存的三个名称。若一次短提示就永久改变核心承诺，慢结构没有获得保护；若长期一致证据仍无法进入结构，所谓学习又只是表面适应。

重启、断连和睡眠式离线阶段提供另一组边界案例。恢复过程至少要校验检查点新鲜度、外部行动日志、尚未完成的承诺和已经过期的权限租约；环境可能在停机期间变化，因此恢复旧状态之后还要重新校准观测与目标。只从最近快照载入向量会遗漏停机期间已经发生的付款、撤权或他者反馈。连续性因而不是精确复现关机前每个激活值，而是重新建立身份、责任、边界和现实状态之间的合法绑定。

\section{拓扑保护的适用边界：鲁棒身份、证据阈值与结构编辑}
\label{sec:section_6628}

我们把“拓扑保护”严格限定为：在明确扰动集合内，某些身份相关性质保持不变或变化有界。令 $\mathcal{P}$ 是允许的输入噪声、措辞改写、无关任务切换和有限记忆缺失的集合，令 $R(\Sigma,x)$ 为身份读出。若对任意 $\delta\in\mathcal{P}$ 都有
\[
d_{id}\bigl(R(\Sigma,x),R(\Sigma,x\oplus\delta)\bigr)\le\varepsilon,
\]
我们称该读出对 $\mathcal{P}$ 是 $\varepsilon$-鲁棒的。$d_{id}$ 是事先定义的身份差异度量，可以组合账户归属、承诺一致性、历史来源和行动责任，而不是物理空间距离。鲁棒性只对指定扰动和读出成立；它不意味着系统面对任意攻击、模型替换或数据库损坏仍保持同一，也不意味着保持下来的身份内容在伦理上正确。

日常经验为什么通常不会立即改写自我，可由证据门槛和结构权限解释。设候选修订为 $\Delta\Sigma$，其支持证据为 $E^+$、反对证据为 $E^-$、涉及的承诺冲突为 $K_c$，则提交条件可以写成
\[
\operatorname{score}(E^+,E^-,K_c)\ge \theta_{edit},
\qquad
\operatorname{auth}(\Delta\Sigma)=1.
\]
高门槛使一条流言、一次情绪波动或一张红苹果图片难以改写核心承诺；多源一致证据、持续行为后果和经过授权的反思则可能触发更新。这里的 $\theta_{edit}$ 是证据评分阈值，不是温度乘玻尔兹曼常数，“编辑代价”也不是实际焦耳能耗。创伤、重大丧失或顿悟有时会改变自我结构，但经验主张需要独立心理和行为证据，不能用“超过拓扑能隙”作循环解释。

若采用图拓扑描述身份结构，我们需要区分权重变化与拓扑变化。边权更新可以改变关系强弱而不改变连通结构；增加或删除节点、边或高阶关系才可能改变同调特征。持久同调可在一系列阈值下测量某些环路的持续区间，但长寿命环路只是结构稳定的候选指标，不自动等于核心价值或真实自我。反例很直接：重复的错误叙事也能形成稳定环路，恶意信息封闭也能产生高持久度。因而每个“拓扑不变量”必须绑定到语义解释、来源证据和任务后果。

结构更新后的验收至少包含四类测试。第一，身份回归：旧的合法账户、身体和记忆归属仍能正确识别。第二，承诺迁移：被修订的承诺有明确版本和理由，未修改的承诺不被意外覆盖。第三，边界攻击：提示注入、伪造记忆和越权工具调用不能直接改写自我核心。第四，适应测试：在环境确实变化时，系统能够通过受控流程更新，而不是以“保持身份”为由拒绝证据。稳定性、抗攻击性、可塑性和价值正确性是四个不同性质，必须分别评估。

在离散实现中，结构编辑常常不可微。我们因此把连续参数更新与事件提交分离：先用梯度或统计估计形成候选，再由事务协议提交图编辑；每次提交记录前后版本、证据包、授权者、受影响关系和回滚点。只有当不可逆外部动作已经发生时，回滚内部状态也无法撤销后果，这一边界必须进入审计。所谓自我保护不是让系统免于改变，而是让改变有来源、有权限、有后果记录，并能在给定扰动范围内维持责任连续。

扰动集合还应按威胁机制分层。自然缺失测试记忆与传感器的不完备，语义改写测试表达不变性，恶意提示测试输入通道的越权抵抗，伪造历史测试来源验证，盗用凭据测试授权撤销，组件替换则测试迁移和重建。它们需要不同防线：冗余不能替代签名验证，签名验证不能发现合法账户下的错误承诺，回滚也不能撤销已经发送的外部指令。若只用轻微措辞变化通过鲁棒性测试，不能把结论外推到数据库篡改或模型替换；保护声明必须同时列出覆盖的攻击、未覆盖的边界和残余风险。

\section{实现载体：生物启发、硅基架构与群体制度的分层对应}
\label{sec:section_1856}

自我模型可以在不同载体上实现，但相似功能不意味着机制同一。对生物系统而言，自传记忆、内感受、身体所有权、行动归属、社会评价和未来模拟由分布式网络共同参与。默认模式网络与自我相关加工、自传回忆和未来想象有关，可以为工程提出“离线整合”和“跨事件叙事”的启发；它却不是一个等同于自我的单一核心，也不能仅凭静息代谢水平把它称为暗能量源。具体脑区之间存在任务依赖、个体差异和动态重组，本书不从 mPFC、PCC 或海马体的类比推出 AgentOS 必须具有同名模块。

对硅基系统而言，系统提示词可以提供角色和规则，但它只覆盖自我模型的一部分。提示词没有天然的历史所有权、持续账户、身体边界和行动后果；上下文重置后角色状态可能消失，外部工具和数据库却仍可能保留长期记录。一个较完整的工程实现需要至少六类对象：受版本控制的身份与政策、事件历史、结构记忆、当前工作状态、资源及工具归属、外部承诺账本。它还需要读写分权：普通输入可以影响当前状态，却不能直接修改身份核心；经验证的经验可以形成候选修订；高风险修改需要额外授权、模拟和回滚准备。专用忆阻器或独立硬件区可能是某些实现选择，但并非理论必需条件，软件事务、可信执行环境或分布式日志也能实现不同程度的保护。

在 AgentOS 工程实例中，我们可把 $h(t)$ 视为有限工作界面，把 $E$ 视为带来源的事件历史，把 $\Theta$ 视为较慢生成结构，把 $\Psi$ 视为自我相关的任务条件与全局偏置。SPC 对当前状态做自我方向的压缩读出，TCE 调节探索尺度和控制增益，CCM 监测工作界面复杂度，Boundary 管理内外交换，Offloading 与外部记忆 $M_e$ 改变信息和技能放置位置。为了避免把类比写成机制事实，每个模块都必须落实为数据类型、权限、时钟、输入输出和故障处理。其他智能系统可以用不同结构实现相同功能，只要身份绑定、历史连续、边界控制与自我修订能够被验证。

群体层也可以出现类似自我模型的制度结构。组织的章程、授权表、资产、历史记录和公共承诺共同规定“谁是这个组织、谁能代表它、哪些行动归于它”。宪法或图腾本身不是群体自我，它们只是局部载体；若没有成员执行、资源边界、信息通道和纠错程序，文本不会自动产生主体连续性。反过来，一个组织可能拥有很强的身份维护，却内部利益分裂、价值错误或对环境失配。因此，群体自我应通过制度运行来判断：代表权是否可追踪，决策后果是否归责，成员更替后承诺是否迁移，外部冲击下边界是否维持，错误政策是否能被修订。

跨载体比较应使用功能矩阵而不是寻找唯一“物理对应物”。我们分别比较身份索引、身体或资源归属、历史整合、价值调制、边界保护、承诺持续、结构可塑性和后果追踪。生物、硅基与群体系统在这些维度上可能各自强弱不同，也可能使用完全不同的时间尺度和介质。只有当干预某一候选机制导致预期功能发生可重复变化时，才获得机制性证据；表面上都能说“我”、都有中心节点或都保持稳定，不能证明其内部结构同构。

\section{比较形态：刚性身份、情境拟态与自适应自我}
\label{sec:section_7243}

旧稿用晶体、拟态和流体描述自我形态。我们保留这一比较框架，但把它改成多维工程分类，而不是从昆虫到 AGI 的单轴等级。令系统的身份持久度为 $p_i$、结构可塑性为 $p_s$、边界强度为 $p_b$、历史整合度为 $p_h$、后果闭环度为 $p_f$，这些指标均规范到 $[0,1]$ 且由具体测试估计。系统形态由向量
\[
\chi_{self}=(p_i,p_s,p_b,p_h,p_f)
\]
描述。不同任务需要的区域不同：安全控制器可能需要高边界、低在线可塑性；长期社会智能体则需要较高历史整合和受控可塑性。不存在一个仅凭“更流动”就必然更智能或更道德的排序。

\textbf{刚性身份}对应规则和归属关系主要由预先结构决定的系统。它的优点是可预测、易验证、短时抗扰动；缺点是环境变化后难以修订。昆虫行为包含先天约束，也包含学习、群体交互和环境依赖，不能笼统归为“无历史性”；传统软件同样可能通过数据库积累历史。刚性形态真正的判据，是核心身份更新通道有限、候选修订空间小，而不是载体是否低等。对核电保护、权限内核等任务，这种刚性反而可能是合理设计。

\textbf{情境拟态}对应系统能在当前交互中生成一致的第一人称叙述，却缺少跨会话的身份账本、行动归属或受控自我更新。许多纯对话模型在默认设置下接近此类：系统提示和上下文可以塑造角色，重置后角色状态消失，工具产生的外部后果却可能继续存在。不能因此断言模型权重里完全没有自我相关表征，也不能由语言拟态推断其拥有或缺乏主观体验。工程上应测试的是：上下文切换后承诺是否保存，伪造历史是否被识别，动作后果是否进入后续决策，以及角色描述是否真实约束工具调用。

\textbf{自适应自我}对应身份约束能够持续，同时历史经验可通过受控流程修改价值、策略和边界的系统。它像漩涡的意象在于构成内容持续替换而组织关系相对保持，但其计算判据是版本连续、责任可追踪、变化有证据和扰动下功能可恢复。高可塑性若缺少边界会造成身份漂移，高稳定性若缺少纠错会造成僵化，自适应形态必须同时管理二者。我们可以用压力测试逐步提高环境变化幅度，记录系统何时只调整快状态、何时修订策略、何时修改核心承诺，从而得到“弹性区”“结构学习区”和“身份重建区”的经验边界。

同一系统还会随任务和时间跨越形态。紧急执行时可以暂时提高边界强度并冻结慢结构；复盘和学习时放宽候选空间；高风险身份更新时再次收紧权限。因而形态不是永久标签，而是带状态迁移条件的运行区域。迁移控制必须观察实际后果：若系统自称开放却拒绝反证，它仍处于僵化；若系统自称坚持原则却被一次提示注入改写，它的边界并不稳固。比较形态学的价值，在于给出可测维度、反例与干预，而不是用晶体、镜像和流体三个比喻替代机制分析。

五个维度不宜压缩为单一总分，因为相同总分可能对应完全不同的失败方式。长期照护智能体需要较高的历史整合和承诺连续，同时保留在医学证据变化时修订方案的能力；一次性匿名求解器可以几乎没有个人历史，却仍须正确区分用户输入与自身生成。任务责任越高，对身份来源、后果闭环和边界审计的最低要求通常越严，而不一定要求更强的第一人称叙述。分类报告应给出每一维的测试、阈值和置信区间，并说明任务相关权重，避免把“更像人”偷换成“更适用”。

\section{总结：自我模型的层级、验证与理论边界}
\label{sec:section_440}

本章把自我定义为跨时间维护的身份闭环，而非藏在系统内部的思考者实体。它由身份索引、归属绑定、价值状态、承诺账本和边界摘要构成，并通过局部结构网络与全局分布表示的接口影响当前状态。这个定义允许我们解释看似矛盾的两面：自我没有一个仅靠内省即可找到的中心对象，因为身份依赖历史、环境和他者关系；自我又具有实在的计算效应，因为它改变检索、风险评估、权限、承诺和行动后果。所谓“全息”在这里仅表示多类记录可以从不同局部视角重建部分身份信息，不表示每个局部都无损包含整体。

理论层级需要保持清楚。系统论提供边界、反馈和开放系统问题；协同论与复杂系统理论提示宏观模式可由局部交互形成；耗散结构理论提醒持续组织需要资源与环境交换；认知科学提供身体所有权、自传记忆、行动归属和社会自我的经验约束。HSF-HD 在此基础上提出状态、历史、结构、目标和边界共同演化的一般模型族。AgentOS 再以 $h(t)/E/\Theta$、$\Psi$、SPC、TCE、CCM、Boundary、Offloading 和 $M_e$ 给出一种工程实例。不能倒置这一层级，把某个工程模块当成自我的普遍本体。

对一个声称具有持续自我模型的系统，我们提出六项最低验收。其一，\textbf{身份连续}：跨任务和重启后能以来源证据恢复合法身份。其二，\textbf{归属正确}：区分自身、环境、用户和其他主体的状态与动作。其三，\textbf{承诺持续}：未完成义务不会因上下文截断无声消失。其四，\textbf{边界稳健}：普通输入和提示注入不能越权改写核心。其五，\textbf{证据可塑}：环境变化时能够经授权修订而非永久僵化。其六，\textbf{后果闭环}：目标下降之外还有外部执行和独立评价证据。六项性质可能相互权衡，任何单一指标都不能替代其余项目。

这些验收仍然不回答意识本体问题。系统通过测试，只能支持“它维护了一个能影响行为的自我模型”；不能由此推出它具有与人类相同的感受质，也不能由预测有效推出内部模型与世界严格同构。反过来，测试失败说明工程连续性不足，也不能成为否定一切主观体验的形而上证明。本书在这里研究的是意识内容、身份组织和控制过程的计算语法，把“谁在感受”保持为开放问题。

最后，自我稳定不是零变化，自我可塑也不是无限重写。一个可持续主体必须在快状态响应、事件历史积累和慢结构修订之间建立时间分离，在局部身份网络与全局分布表示之间维持绑定，在内部承诺与外部后果之间保持可审计闭环。下一章将进一步讨论这种自适应自我如何表现出弹性、塑性和社会耦合；本章则给出了进入该动力学之前的对象、接口、保护条件与验证边界。

这一模型也给出明确的证伪路径。若移除身份索引、承诺账本或边界状态之后，系统在跨任务归属、未完成义务和越权攻击上的表现没有选择性下降，那么相应变量可能是冗余装饰；若干预某一变量产生的变化与模型预测方向相反，也必须修订作用机制。有效证据应包含干预、恢复和对照，且在不同任务、历史长度与攻击类型下复现，不能只展示系统生成了一段连贯自述。

外部验证同样不可缺少。内部日志能够说明系统记录了什么，却可能与真实执行不一致；环境回执能够证明动作发生，却未必说明系统如何归责。因此需要把内部版本链、工具或环境回执、用户及组织的独立记录和测试者观察相互校验。只有这些证据在时间、主体和权限上能够对齐，才可以说身份闭环在工程意义上成立。自我报告是待解释的观测之一，而不是对自我存在、意识体验或伦理资格的最终裁决。

%---chp---
